\documentclass[11pt]{article}

% ============================ packages ============================
\usepackage[utf8]{inputenc}
\usepackage[T1]{fontenc}
\usepackage[a4paper,margin=2.2cm]{geometry}
\usepackage{amsmath,amssymb,amsthm,mathtools}
\usepackage{xcolor}
\usepackage{booktabs}
\usepackage{enumitem}
\usepackage{listings}
\usepackage{tcolorbox}
\usepackage{graphicx}
\usepackage[hidelinks]{hyperref}
\usepackage{titlesec}
\tcbuselibrary{skins,breakable}

% ============================ colours ============================
\definecolor{accent}{HTML}{1F6FB2}
\definecolor{accent2}{HTML}{2CA25F}
\definecolor{warn}{HTML}{C0392B}
\definecolor{mathc}{HTML}{6C5CE7}
\definecolor{codebg}{HTML}{F4F6F8}
\definecolor{asciibg}{HTML}{FBFAF5}
\definecolor{codekw}{HTML}{8E44AD}
\definecolor{codestr}{HTML}{27AE60}
\definecolor{codecom}{HTML}{7F8C8D}

% ============================ code style ============================
\lstdefinestyle{py}{
  language=Python, basicstyle=\ttfamily\footnotesize,
  backgroundcolor=\color{codebg}, keywordstyle=\color{codekw}\bfseries,
  stringstyle=\color{codestr}, commentstyle=\color{codecom}\itshape,
  showstringspaces=false, breaklines=true, frame=single, framerule=0pt,
  rulecolor=\color{codebg}, xleftmargin=6pt, xrightmargin=6pt, aboveskip=6pt, belowskip=6pt,
  literate={~}{{$\sim$}}1
}
\lstdefinestyle{ascii}{
  basicstyle=\ttfamily\small, backgroundcolor=\color{asciibg},
  frame=single, framerule=0pt, rulecolor=\color{asciibg},
  xleftmargin=6pt, xrightmargin=6pt, aboveskip=6pt, belowskip=6pt, columns=fullflexible
}
\newcommand{\code}[1]{\texttt{\detokenize{#1}}}

% ============================ boxes ============================
\newtcolorbox{plainbox}[1][In plain words]{colback=accent!6,colframe=accent,
  fonttitle=\bfseries,title=#1,breakable,boxrule=0.6pt,left=4pt,right=4pt,top=3pt,bottom=3pt}
\newtcolorbox{mathbox}[1][If you want the formula]{colback=mathc!5,colframe=mathc,
  fonttitle=\bfseries,title=#1,breakable,boxrule=0.5pt,left=4pt,right=4pt,top=3pt,bottom=3pt}
\newtcolorbox{tipbox}[1][Exam tip]{colback=accent2!7,colframe=accent2,
  fonttitle=\bfseries,title=#1,breakable,boxrule=0.6pt,left=4pt,right=4pt,top=3pt,bottom=3pt}
\newtcolorbox{pitbox}[1][Easy to mix up]{colback=warn!6,colframe=warn,
  fonttitle=\bfseries,title=#1,breakable,boxrule=0.6pt,left=4pt,right=4pt,top=3pt,bottom=3pt}
\newtcolorbox{churnbox}{colback=black!4,colframe=black!55,
  fonttitle=\bfseries,title=How this maps to your FictiPay churn project,breakable,
  boxrule=0.6pt,left=4pt,right=4pt,top=3pt,bottom=3pt}

\titleformat{\section}{\Large\bfseries\color{accent}}{\thesection}{0.6em}{}
\titleformat{\subsection}{\large\bfseries}{\thesubsection}{0.6em}{}

\setlength{\parindent}{0pt}
\setlength{\parskip}{4pt}

% ============================ title ============================
\title{\vspace{-1.5cm}\textbf{Survival Analysis --- The Friendly Guide}\\[3pt]
\large Intuition first, formulas optional \\[2pt]
\normalsize NSUCEC Datathon · Team Cybernauts}
\author{}
\date{}

\begin{document}
\maketitle
\vspace{-1.0cm}
\hrule
\vspace{6pt}

\begin{plainbox}[How to read this guide]
Every topic starts with a \textbf{plain-English story} (usually about customers leaving an app, or
light bulbs burning out). The \textbf{\textcolor{mathc}{purple boxes}} hold the actual formula --- read
them only if you want the math; you can understand everything without them. The grey boxes connect
each idea back to the churn project you already built.
\end{plainbox}

\tableofcontents
\vspace{6pt}
\hrule

% =================================================================
\section{The One Big Idea}
% =================================================================

Survival analysis answers a single question: \textbf{``how long until something happens?''} --- where
the ``something'' (called the \textbf{event}) is usually bad: a person dies, a machine breaks, a
customer churns.

Here is the twist that makes it special. Imagine you switch on \textbf{100 light bulbs} and watch
them. Some burn out and you write down exactly when. But after 6 months you stop watching, and
\textbf{30 bulbs are still glowing}. You do \emph{not} know how long those will last --- only that
it is \emph{more than 6 months}.

The naive move is to throw those 30 away (``incomplete data''). That is a terrible idea ---
\emph{those are your best bulbs!} Ignoring them makes your product look worse than it really is.

\begin{plainbox}[The whole point of the field]
Survival analysis is simply \textbf{the correct way to use ``still alive when I stopped watching''
information.} That half-known data point has a name --- \textbf{censoring} --- and handling it
properly is the entire reason this toolbox exists.
\end{plainbox}

\begin{churnbox}
A customer = a light bulb. Churning = burning out. A customer who is \emph{still active} at your
April 1 cutoff = \textbf{a bulb still glowing = censored}. You know they ``survived'' (kept
transacting) at least that long --- you just do not know \emph{when} they will finally quit.
\end{churnbox}

% =================================================================
\section{The Two Ways to Look at Survival}
% =================================================================

There are two questions you can ask about your bulbs (or customers), and they are just two views of
the \emph{same} story.

\subsection{The survival curve $S(t)$ --- ``what \% are still alive?''}
Starts at 100\%, only ever goes \emph{down}. At month 3, maybe 80\% are still glowing; at month 6,
60\%. That downward curve is the \textbf{survival function}. Simple.

\subsection{The hazard $h(t)$ --- ``how dangerous is right now?''}
This is the subtle one. The hazard is the \textbf{risk rate for the survivors only}: of those still
alive \emph{at this very moment}, what is the chance they fail in the next instant?

A bulb that has already lasted 5 years might be \emph{more} likely to pop soon --- it is worn out.
A brand-new customer might be fragile too (never really got hooked). The hazard captures ``how risky
is \emph{this particular moment}, given you have made it this far.''

\begin{plainbox}[The mental picture]
Think of the \textbf{hazard} as a \emph{danger meter} ticking moment by moment, and the
\textbf{survival curve} as the \emph{scoreboard}. The danger meter is what \emph{drives} the
scoreboard down. High hazard now $\Rightarrow$ the survival curve is about to drop steeply.
\end{plainbox}

\begin{mathbox}
Four ways to describe the same thing, all interchangeable:
$$S(t)=P(T>t)\ \ (\text{still alive}),\qquad h(t)=\frac{f(t)}{S(t)}\ \ (\text{risk rate now}),$$
$$H(t)=\int_0^t h(u)\,du\ \ (\text{risk piled up so far}),\qquad \boxed{S(t)=e^{-H(t)}}.$$
That last boxed identity is the bridge: total accumulated danger $H$ determines how many survive.
\end{mathbox}

\begin{churnbox}
Your best features --- \code{recency_days}, \code{gap_z}, \code{p_silent30_poisson} --- are
basically \textbf{hazard readings}. A customer silent for 40 days when they normally transact daily
has a \emph{sky-high} hazard right now (danger meter pegged). That is exactly why those features
dominated your model.
\end{churnbox}

% =================================================================
\section{Kaplan--Meier: Drawing the Curve from Messy Data}
% =================================================================

\textbf{Kaplan--Meier (KM)} is just the method for \emph{drawing} that ``\% still alive over time''
curve when some data is censored. You do not need the formula --- you need the picture.

\begin{lstlisting}[style=ascii]
 % still alive
 100% |#####
      |    #####            <- flat stretch: nobody churned here
  80% |        ###             (some may be CENSORED here -- they
      |           ####          just leave the pool quietly,
  60% |              ######     they do NOT cause a drop)
      |                   ###
  40% |                      ##  <- each STEP DOWN = one real churn
      +---------------------------> time
\end{lstlisting}

It is a \textbf{staircase going down}. Every step down = one real event (a churn). Flat stretches =
nobody churned. The clever part: when someone is \textbf{censored} (still active when you stop
watching), they do \emph{not} cause a step down --- they just \textbf{quietly leave the pool} so
they are not counted in later steps. That is how KM honestly uses your ``still-glowing bulbs.''

\begin{plainbox}[Two things people always ask about]
\textbf{Median survival} = the time where the staircase crosses 50\% (the ``typical'' lifespan). If
the curve never reaches 50\% (lots of censoring), the median is simply ``not reached'' --- a
perfectly valid answer. \textbf{Why no drop at a censoring time?} Because a censored person never had
the event; they only stopped being watched.
\end{plainbox}

\begin{mathbox}
At each event time $t_i$: $d_i$ = how many failed, $n_i$ = how many were still at risk just before.
$$\hat S(t)=\prod_{t_i\le t}\Big(1-\frac{d_i}{n_i}\Big).$$
Each bracket is ``fraction who survived this instant''; multiply them along the staircase.
\end{mathbox}

\subsection{Comparing two groups: the log-rank test}
Say you split customers into \textbf{high-balance vs low-balance}, draw a KM staircase for each, and
the high-balance one stays higher (they churn less). Is that a \emph{real} difference, or just luck?

\begin{plainbox}
The \textbf{log-rank test} answers exactly that. It tallies ``how many churned in each group, and
when'' versus ``what we would expect if the two groups were secretly identical.'' A big mismatch
$\Rightarrow$ the groups are genuinely different. It gives you a \textbf{yes/no with a confidence
level (p-value)} --- but \emph{not} ``by how much.'' For ``by how much,'' you need the Cox model
(next).
\end{plainbox}

\begin{pitbox}
Kaplan--Meier and the log-rank test only compare \textbf{groups} (categories). The moment you want
to handle a continuous number (balance, age, \code{gap_z}) or several factors at once, you must move
to a \emph{regression} model --- Cox or Weibull.
\end{pitbox}

% =================================================================
\section{The Cox Model: Many Factors at Once}
% =================================================================

KM compares groups. But real customers have \textbf{lots} of traits at the same time --- age,
balance, region, recency, transaction count. You cannot draw a separate curve for every combination.

The \textbf{Cox proportional hazards model} fixes this. Its whole idea in one sentence:

\begin{plainbox}[The Cox idea]
\emph{Everyone shares one baseline ``danger-over-time'' shape, and your personal traits
\textbf{multiply} that baseline up or down.}
\end{plainbox}

The output for each trait is a \textbf{hazard ratio (HR)} --- the multiplier:
\begin{itemize}[leftmargin=1.4em]
  \item \textbf{HR $=2$} $\Rightarrow$ that trait \emph{doubles} the churn rate (a risk factor).
  \item \textbf{HR $=0.5$} $\Rightarrow$ \emph{halves} it (protective).
  \item \textbf{HR $=1$} $\Rightarrow$ no effect.
\end{itemize}

Example: ``each extra week of silence has HR $=1.4$'' means every quiet week \textbf{multiplies}
churn risk by $1.4\times$. That is a clean, boardroom-friendly sentence --- far more useful than a
raw probability.

\begin{mathbox}
$$h(t\mid \mathbf x)=\underbrace{h_0(t)}_{\text{shared baseline}}\times
\underbrace{e^{\beta_1 x_1+\dots+\beta_p x_p}}_{\text{your personal multiplier}},
\qquad \mathrm{HR}_j=e^{\beta_j}.$$
``Semi-parametric'' = it never has to guess the baseline shape $h_0(t)$; that cancels out in the math.
\end{mathbox}

\subsection{The one catch: the ``proportional hazards'' assumption}
\begin{plainbox}
Cox \textbf{assumes your multiplier stays the same over time}. If silence always multiplies risk by
$1.4\times$ --- whether the customer is brand new or two years old --- great, Cox is valid. But if a
trait matters a lot \emph{early} and fades \emph{later}, Cox is wrong and you must check it.
\end{plainbox}

\textbf{How to check it (know one or two):}
\begin{itemize}[leftmargin=1.4em]
  \item \textbf{The parallel-lines test} (log--log plot): transform each group's curve a certain way
  and look. \textbf{Stay parallel} $\Rightarrow$ assumption holds. \textbf{Cross or fan apart}
  $\Rightarrow$ broken.
  \item \textbf{Schoenfeld residuals}: a built-in statistical test (\code{cph.check_assumptions} in
  \code{lifelines}) that flags any trait whose effect drifts with time.
\end{itemize}

\textbf{If it is broken:} split that trait into separate baselines (``stratify''), let its effect
change over time, or switch to a Weibull / random-survival-forest model that does not require the
assumption.

\begin{tipbox}
Cox tells you \emph{relative} risk (``$1.4\times$ more likely''), not \emph{when}. If the question is
``\emph{how many days sooner} will they churn?'', that is the parametric / AFT model's job.
\end{tipbox}

% =================================================================
\section{Weibull: Assuming a Shape for the Risk}
% =================================================================

Cox leaves the baseline danger-meter shape a \emph{mystery} (just traced from the data). A
\textbf{parametric} model instead \emph{commits} to a known mathematical shape. \textbf{Weibull} is
the famous flexible one, and it has a single \textbf{``shape knob,'' called $k$}, that is the
most-tested fact in this whole topic:

\begin{plainbox}[The Weibull shape knob $k$]
\begin{itemize}[leftmargin=1.2em,topsep=2pt]
  \item \textbf{$k>1$ --- danger keeps \emph{rising}.} Things wear out (old machines, aging bulbs).
  \item \textbf{$k=1$ --- danger is \emph{flat}.} Same risk forever, no matter how long it has been
  (purely random failures). This is the plain ``exponential'' case.
  \item \textbf{$k<1$ --- danger \emph{falls} over time.} Survive the shaky beginning and you are
  golden. \textbf{This is classic for customers:} many churn in the first few weeks, but whoever
  sticks around becomes loyal for the long haul.
\end{itemize}
\end{plainbox}

\textbf{Why use Weibull instead of Cox?} Three reasons:
\begin{enumerate}[leftmargin=1.5em]
  \item \textbf{Predict the future.} Cox cannot say anything past your last observed time; a Weibull
  curve keeps going, so you can forecast.
  \item \textbf{Smooth curves} and the ability to \emph{simulate} customer lifetimes.
  \item \textbf{The shape itself is the insight} (``our churn is $k<1$ $\Rightarrow$ fix onboarding,
  that is where we bleed people'').
\end{enumerate}

\begin{mathbox}
$$S(t)=e^{-(t/\lambda)^k},\qquad h(t)=\frac{k}{\lambda}\Big(\frac{t}{\lambda}\Big)^{k-1}.$$
$k$ = shape (rising/flat/falling), $\lambda$ = scale (overall ``stretch'' of the timeline).
The related \textbf{AFT} view writes time directly: a trait's ``time ratio'' $e^{\beta}$ says the
event happens that many times \emph{sooner} or \emph{later}.
\end{mathbox}

% =================================================================
\section{The Machine-Learning Versions}
% =================================================================

When relationships get complicated (lots of interacting features, no clean assumptions), you bring
in machine learning --- the \emph{same} model families you already used, just taught to handle
\emph{time + censoring}.

\begin{itemize}[leftmargin=1.4em]
  \item \textbf{Random Survival Forest (RSF)} --- a random forest where each tree predicts
  time-to-event and knows how to handle censored customers. It splits the data to separate
  ``high-risk soon'' from ``low-risk'' groups as sharply as possible. No proportional-hazards
  assumption; finds non-linear patterns automatically.
  \item \textbf{XGBoost AFT} --- gradient boosting (your XGBoost / CatBoost world) that predicts the
  \emph{actual time until churn} and properly uses ``still active'' rows. You feed each customer a
  \emph{range} for their event time: for a still-active customer, ``somewhere after today''
  (lower bound = today, upper bound = infinity).
\end{itemize}

\begin{plainbox}
Same machinery you already know from the churn round. The only new trick: these output a
\textbf{``when''} (a time or a survival curve), not just a yes/no.
\end{plainbox}

\begin{churnbox}
You solved churn as a \emph{yes/no at a fixed date} (``any transaction in April?''), which was the
right call because the target is defined on a fixed window. But underneath, it is a survival problem:
time-to-next-transaction, censored at the cutoff. If a judge asks \emph{``why not survival
analysis?''}, the crisp answer is: \emph{``our label is a fixed-time yes/no, so classification
directly optimises the scored metric. Survival would be the right tool if we needed churn
\textbf{timing} or a churn-by-any-future-date curve.''}
\end{churnbox}

% =================================================================
\section{Grading a Survival Model}
% =================================================================

\subsection{C-index --- ``does it rank people correctly?''}
\begin{plainbox}
The \textbf{C-index} is the survival version of \textbf{AUC}. Take any two customers: did the model
correctly guess which one churns sooner? The fraction it gets right is the C-index.
\textbf{0.5 = coin flip, 1.0 = perfect.} It measures \emph{ranking / ordering}.
\end{plainbox}

\subsection{Brier score --- ``are the probabilities accurate?''}
\begin{plainbox}
The \textbf{Brier score} checks whether the actual predicted \emph{probabilities} are right, not just
the ranking. \textbf{Lower = better.} It catches a model that ranks customers fine but outputs silly
probability numbers (e.g.\ says ``90\% sure'' and is wrong half the time).
\end{plainbox}

\begin{tipbox}
Report \textbf{both}: a model can rank well (good C-index) but be badly calibrated (bad Brier), or
the reverse. C-index = \emph{discrimination}; Brier = \emph{calibration $+$ discrimination}.
This contrast is a favourite exam question.
\end{tipbox}

% =================================================================
\section{The Two Python Toolboxes}
% =================================================================

\begin{itemize}[leftmargin=1.4em]
  \item \textbf{\code{lifelines}} --- the friendly one. Best for Kaplan--Meier curves, the Cox model,
  and clear interpretation/plots.
  \item \textbf{\code{scikit-survival} (\code{sksurv})} --- sklearn-style. Best for the ML models
  (RSF) and the scoring metrics (C-index, Brier).
\end{itemize}

\textbf{The bare-minimum code you should recognise:}
\begin{lstlisting}[style=py]
# ---- lifelines: Kaplan-Meier curve ----
from lifelines import KaplanMeierFitter, CoxPHFitter
kmf = KaplanMeierFitter()
kmf.fit(durations=df["time"], event_observed=df["event"])
kmf.plot_survival_function()       # the staircase
kmf.median_survival_time_          # the typical lifespan

# ---- lifelines: compare two groups ----
from lifelines.statistics import logrank_test
res = logrank_test(a["time"], b["time"], a["event"], b["event"])
print(res.p_value)                 # small p  ->  groups really differ

# ---- lifelines: Cox model (hazard ratios) ----
cph = CoxPHFitter().fit(df, duration_col="time", event_col="event")
cph.print_summary()                # look at exp(coef) = the HR
cph.check_assumptions(df)          # is the PH assumption ok?

# ---- scikit-survival: Random Survival Forest + scoring ----
from sksurv.util import Surv
from sksurv.ensemble import RandomSurvivalForest
from sksurv.metrics import concordance_index_censored
y = Surv.from_arrays(event=df["event"].astype(bool), time=df["time"])
rsf = RandomSurvivalForest(n_estimators=300, n_jobs=-1).fit(X, y)
risk = rsf.predict(X_test)
c = concordance_index_censored(y_test["event"], y_test["time"], risk)[0]
\end{lstlisting}

% =================================================================
\section{30-Second Recap (read this the morning of)}
% =================================================================
\begin{plainbox}
Watch a bunch of customers. Some churn (you see \emph{when}); some are \textbf{still active when you
stop watching} (\textbf{censored} --- the whole point of the field). \textbf{Kaplan--Meier} draws the
survival staircase. \textbf{Log-rank} tests whether two groups' staircases really differ.
\textbf{Cox} says how each trait \emph{multiplies} churn risk (the \textbf{hazard ratio}).
\textbf{Weibull} assumes a shape and tells you whether risk \emph{rises, stays flat, or falls} over a
customer's life (the shape knob $k$). \textbf{RSF / XGBoost-AFT} are the machine-learning versions.
\textbf{C-index} grades the ranking; \textbf{Brier} grades the probabilities.
\end{plainbox}

% =================================================================
\section{One-Page Cheat Sheet}
% =================================================================
\begin{center}
\renewcommand{\arraystretch}{1.35}
\begin{tabular}{@{}p{3.1cm}p{12.3cm}@{}}
\toprule
\textbf{Term} & \textbf{What it means in one line} \\
\midrule
Event & The thing you are timing (death, failure, churn). \\
Censoring & ``Still alive when we stopped watching'' --- known partly, not fully. The reason this field exists. \\
Right-censored & Most common: event has not happened yet at the cutoff. \\
Left truncation & Subject only enters the study late; you could not have seen earlier failures. \\
$S(t)$ survival & \% still alive at time $t$ (the scoreboard; only goes down). \\
$h(t)$ hazard & Risk rate \emph{right now}, for survivors only (the danger meter). \\
Kaplan--Meier & The downward staircase; steps at churns, censored people leave quietly. \\
Median survival & Where the staircase crosses 50\% (``typical'' lifespan); may be ``not reached''. \\
Log-rank test & Do two groups' curves really differ? Gives a yes/no (p-value), not a size. \\
Cox model & Each trait \emph{multiplies} a shared baseline risk. \\
Hazard ratio (HR) & The multiplier: $>1$ riskier, $<1$ protective, $=1$ no effect. \\
PH assumption & Cox assumes the multiplier stays constant over time; check with parallel-curve / Schoenfeld tests. \\
Weibull shape $k$ & $k>1$ risk rises, $k=1$ flat, $k<1$ risk falls (early churn then loyal). \\
Parametric vs Cox & Use parametric to forecast the future / get smooth curves / read the shape. \\
RSF & Random forest for time-to-event; no PH assumption; finds interactions. \\
XGBoost AFT & Boosting that predicts the actual time; censored rows get a time \emph{range}. \\
C-index & ``Did it rank who-fails-first correctly?'' = survival AUC (0.5 random, 1.0 perfect). \\
Brier score & ``Are the predicted probabilities accurate?'' Lower is better. \\
lifelines & Friendly library: KM, Cox, log-rank, plots. \\
scikit-survival & sklearn-style: RSF, C-index, Brier. \\
\bottomrule
\end{tabular}
\end{center}

% =================================================================
\section{Likely Questions (with plain answers)}
% =================================================================
\begin{enumerate}[leftmargin=1.6em]
  \item \emph{Why not just use normal regression?} --- Because of censored people: deleting them
  biases you, and pretending they churned at the cutoff is a lie. Survival math uses them correctly.
  \item \emph{Hazard vs probability?} --- Hazard is a \emph{rate} for those still alive right now; it
  can be bigger than 1 and is not itself a probability.
  \item \emph{Why doesn't the KM curve drop when someone is censored?} --- They never churned; they
  just left the pool, lowering the ``at-risk'' count for later steps.
  \item \emph{What does HR $=0.7$ mean?} --- 30\% \emph{lower} churn rate at every time point
  (protective) --- assuming the proportional-hazards assumption holds.
  \item \emph{How do you check that assumption, and what if it fails?} --- Parallel-curve (log--log)
  plot or Schoenfeld test; fix by stratifying, letting the effect vary with time, or using
  Weibull / RSF.
  \item \emph{When is Weibull better than Cox?} --- When you need to forecast beyond your data, want
  smooth curves, or the rising/flat/falling \emph{shape} is itself the answer.
  \item \emph{What does $k<1$ tell a business?} --- Risk falls over time: heavy early churn, then
  survivors stay loyal $\Rightarrow$ invest in onboarding.
  \item \emph{C-index vs Brier?} --- C-index = ranking (AUC-like); Brier = are the probabilities
  accurate. Report both.
  \item \emph{Why did your team use classification, not survival, for churn?} --- The label is a
  fixed-window yes/no, so classification directly optimises the scored metric; survival would be for
  predicting churn \emph{timing}.
\end{enumerate}

\vfill
\begin{center}\small\itshape
Compile with \code{pdflatex survival_analysis_guide.tex} (run it twice so the contents page fills in).
\end{center}

\end{document}
